\section{Complete verifier and certificate contract}
\label{sec:implementation}

The primary module \texttt{code/checker.py} is the entire trusted
acceptance program. It imports only Python's standard \texttt{json} and
\texttt{sys} modules. It reads one JSON object, validates the closed goal,
reconstructs the proof conclusion, and accepts exactly when they agree.
Acceptance is exit status zero and the line \texttt{ACCEPT}. Rejection
has nonzero exit status. Python's resource limits can reject sufficiently
large valid inputs; no resource-bound completeness claim is made.

\subsection{Concrete syntax}

The public object has exactly two fields, \texttt{goal} and
\texttt{proof}. All term and formula trees are explicitly checked:
\begin{center}
\begin{tabular}{lll}
\toprule
Mathematical object & JSON form & Binding\\
\midrule
Zero & \texttt{0} & None\\
Successor & \texttt{['s', a]} & None\\
Iterator & \texttt{['i', n, a, h]} & One variable in \texttt{h}\\
Equation & \texttt{['=', a, b]} & None\\
Paired formula & \texttt{['d', B]} & Two variables in \texttt{B}\\
\bottomrule
\end{tabular}
\end{center}
The display uses single quotes for readability; JSON files use double
quotes. A positive integer denotes a bound-variable index, with
index \(1\) the innermost variable. In a \(D\) matrix, index \(1\)
denotes \(y\) and index \(2\) denotes \(x\). Negative integers are
rigid free-variable levels used only inside local proof scopes.
At free scope \(k\) and binder depth \(d\), an integer term must satisfy
\(-k\leq a\leq d\). The number zero always denotes zero, not a variable.
Python booleans are explicitly excluded by exact integer-type tests.

Canonical term validation forbids an iterator with literal zero
counter, recursively everywhere. The shared opening operation rebuilds
every visited iterator using the canonical constructor, so exposing a
zero counter during substitution immediately contracts it. The test is
on the rebuilt counter, not on its original raw syntax.

The proof constructors are as follows:
\begin{center}
\begin{tabular}{ll}
\toprule
Rule & Certificate form\\
\midrule
Use & \texttt{['u', clause-index, witness, antecedent-proof]}\\
Refl & \texttt{['r', term]}\\
Step & \texttt{['x', major-proof, minor-proof]}\\
DE & \texttt{['e', goal, paired-proof, branch-proof]}\\
DI & \texttt{['i', paired-goal, witness, branch-proof]}\\
\bottomrule
\end{tabular}
\end{center}

The context is an internal list of antecedent/body pairs; it is not
an unchecked public environment. A DE branch opens its witness at
the next fresh negative level. A DI branch uses the next fresh level
as its induction/case parameter. The conclusion and chosen witness are
validated in the old scope before either branch opens, which prevents
an eigenvariable from escaping. Ordinary local assumptions are stored
as clauses with antecedent \(0=0\).

\subsection{Why the opening operation needs no shift}

The use of negative levels is an implementation convention, not an
extra axiom. Every substitution image is locally closed: all its
positive indices are bound within that image itself, and all its
external variables have negative labels. Inserting such an image
under an additional binder therefore cannot capture an external
variable. The opening traversal increases its depth under the
iterator body and by two under \(D\), replaces the designated bound
index, and decrements larger indices. The substituted image is not
recursively opened a second time.

This argument would fail for arbitrary open positive-index images.
The validator prevents that input. The proof-reconstruction invariant is
that every returned formula is canonical and well-scoped in the current
scope, and every retained clause has a well-scoped antecedent and a
body with exactly one additional permitted bound variable. Each
constructor preserves that invariant. In particular, Step's
premises are explicitly matched as equations with the required
successor-iterator source.

\subsection{Measurements}

All source measurements count the entire acceptance module.
For comparison with the supplied repository, module, class, and
function docstrings are removed from the AST, then the source is
normalized by CPython 3.12.14 \texttt{ast.unparse}, followed by one LF.
Significant tokens exclude comments, whitespace, indentation, newline,
encoding, and end markers. AST nodes are counted after that same
docstring removal. Gzip uses maximum level \(9\) and timestamp zero;
it is informational, not an optimization objective.
For the delivered primary Python file itself, including its documentation
and original formatting, the corresponding figures are
\StepSourceBytes\ source bytes and \StepSourceGzip\ gzip-9 bytes.
Both raw and normalized compressed files are included.

@@METRICS_TABLE@@

The matched canonical control shares the terms, normalization,
opening operation, formula grammar, public fields, validation, clause
protocol, and logical rules with the new checker. Its arithmetic rules
are successor computation and one-template crossed transport.
Consequently that comparison isolates the arithmetic replacement.
The raw Eval control retains raw zero-counter iterators and the three
arithmetic Eval alternatives. Its comparison also includes the cost
of adding canonical syntax.

The secondary modules in \texttt{code/native/} keep the more general
explicit-sequent Eq/\(Q\)/\(D\) interface with a variable count and
public hypotheses. There the matched normalized control measures
3,118 bytes, 1,010 tokens, and 1,080 AST nodes; Refl/Step measures
3,029 bytes, 981 tokens, and 1,051 AST nodes. The pinned raw Eval
checker measures 3,039 bytes, 959 tokens, and 1,026 AST nodes.
Thus the change has a clear matched win, but its comparison with
raw Eval in that second interface is mixed. These modules are
supplementary evidence, not additional hidden dependencies of
the primary verifier.

\subsection{Verification evidence and its limits}

@@TEST_SUMMARY@@

The explicit-sequent suite supplies actual finite proof translations
in both directions between normalized computation/transport and
Refl/Step. The translators are untrusted producers: each output is
checked by the target's complete verifier. They cover every primitive
constructor of those calculi, and preserve the public variable count,
context, and goal.

The tests specifically address the difficult implementation boundaries:
an induction hypothesis that is actually used, a closure actually
used after \(D\)-elimination, nested term binders, normalization
exposed by substitution, arbitrary equation assumptions, malformed
indices, booleans, forged predecessor sources, and escaping
eigenvariables. An independently written numerical evaluator also
checks the opening operation on finite instances. Numerical tests
are supplementary and are not used to infer the nonstandard-model
interpretation theorem.

The script \texttt{reproduce.py} reruns the validation suites and
recomputes all source metrics and deterministic gzip files.
Its optional \texttt{--pdf} switch rebuilds the standalone
\texttt{report.tex} and compiles the PDF. The inherited reference
files are byte-authenticated against Git blob identifiers at the
frozen commit; \texttt{evidence/provenance.json} records that check.

\section{Full first-order interpretation and research boundary}
\label{sec:interpretation-summary}

The interpretation into ordinary logic uses the entire numerical
domain and interprets each native equality as actual equality.
The \(D\) binder and clauses have precisely the ordinary readings
specified above. Genuine ordinary falsity is interpreted as \(0=S0\).
The ordinary scalar realization explicitly includes the separation
axiom \((0=S0)\to\bot\); positive numerical equations alone do not
supply this ordinary falsity bridge.

For transparency, the complete inherited logical macro construction
is reproduced in Appendix~\ref{sec:choice-introduction}.
It derives implication, conjunction, disjunction, both quantifiers,
and explosion using only selectors, finite iteration, and the
native rules. It also proves the paired-formula round trip
\[
 D_{x,y}B
 \quad\dashv\vdash\quad
 \exists_*x\bigl(B(x,0)\Rightarrow_*\forall_*y\,B(x,y)\bigr).
\]
This is why the result concerns full intuitionistic first-order
arithmetic rather than just an equational or negative fragment.
No excluded middle for quantified formulas is introduced.

The arithmetic bridge is also internal. Its order matters:
first construct predecessor by scalar iteration and equation
induction; then prove faithful equality reflection; then obtain
quantifier-free induction; only then construct pairing and compile
general primitive recursion. The scalar update has one bound
state variable and no hidden recursion-counter argument.

For a target recursion
\[
 f(0,\vec p)=g(\vec p),\qquad
 f(Sn,\vec p)=h(n,f(n,\vec p),\vec p),
\]
the inherited pairing \(\pi\) and projections \(L,H\) give
\[
 V(n,\vec p)=I\bigl(n,\pi(0,g^*(\vec p));
     u.\pi(SL(u),h^*(L(u),H(u),\vec p))\bigr),
 \qquad f^*(n,\vec p)=H(V(n,\vec p)).
\]
Equation induction proves \(L(V(n,\vec p))=n\), and the projection
equations prove both defining equations for \(f^*\).
Each fixed finite primitive-recursive description is compiled
externally. There is no internally total universal evaluator over
arbitrary descriptions.

In the other direction, each fixed scalar iterator is an ordinary
primitive recursion ignoring its counter in the update.
Equation induction proves uniqueness for these fixed recursions,
so an external induction over descriptions proves both internal
numerical round trips. Normalization preserves numerical values,
and the clause translation does not change numerical terms.
This establishes the identity graph \(x=y\) on the whole domain
in both composites, including nonstandard elements.

\subsection{What remains open}

The five-rule inventory retains a syntax-directed case/progress
choice inside DI. It has not removed the substantive distinction
between numerical coverage and equation induction. Further merging
the logical rules would need to preserve that distinction and its
eigenvariable conditions.

The result also exchanges a larger equality interface for a modest
increase in some proof sizes: derived equality composition expands
to two Step nodes and an auxiliary reflexivity. Checker source
size, certificate length, and checking time are separate objectives.
The implementation measurements do not establish an optimum for
any of them.

The next useful formal-verification target is a deeply embedded
Step calculus and its equivalence with
\texttt{ConciseIPRA.NativeZeroNormal}. The present package provides
the necessary derivations and tested finite translations, but
does not label an uncompiled Lean sketch as a verified result.

\clearpage
\appendix
\section{Inherited proof material and attribution}

The next five appendices reproduce the exact mathematical source
fragments from \path{Research/certified_zero/reference/} at the
frozen commit. The comparison script verifies their Git blob hashes.
They are included under the repository's MIT-0 license, retained
as \path{reference/LICENSE-ProveIt.txt}. Their descriptions of
earlier contributions are relative to their original reports.
They do not assert priority for this package.

These fragments use the raw scalar notation and their own earlier
native rule presentation. Sections~\ref{sec:uniform-step-theory}
through~\ref{sec:uniform-step-equivalence} prove the equivalence that
allows their arithmetic bridge to be used here. In particular,
the fragments' full formula-transport rule is a recovered rule,
not an additional primitive of the five-rule calculus.

